\section{Saturation: mean, convolution, and Fourier support}
\label{sec:saturation}

The optimizer in \eqref{eq:one-level-optimizer} is always bounded by the
spectral mean $\bar c=\log\rho(C(a))$.  We now characterize equality without
relaxing the integer composition.

\begin{lemma}[Mean preservation and invertibility]
\label{lem:H-invertible}
The circulant map $x\mapsto H(x)$ defined in \eqref{eq:H-definition}
preserves the phase mean and is invertible on $\R^p$.  Consequently,
\begin{equation}
 H(x)\text{ is constant}\quad\Longleftrightarrow\quad x\text{ is constant}.
 \label{eq:H-constant-iff}
\end{equation}
\end{lemma}

\begin{proof}
Because the kernel weights sum to one, cyclic reindexing gives
\[
 \frac1p\sum_jH_j(x)=\frac1p\sum_jx_j.
\]
For invertibility, evaluate the unnormalized kernel polynomial at a $p$th
root of unity $z$:
\begin{equation}
 \sum_{t=0}^{p-1}d^{p-1-t}z^t
 =\frac{d^p-z^p}{d-z}=\frac{d^p-1}{d-z}.
 \label{eq:H-multiplier}
\end{equation}
It is nonzero since $d\geq2$ and $|z|=1$.  Thus every Fourier multiplier of
$H$ is nonzero, which proves invertibility.  Constants are preserved, so
\eqref{eq:H-constant-iff} follows.
\end{proof}

For later use, let $N\geq1$, $m\in\Comp_p(N)$, and set
\begin{equation}
 \Psi_j^{(N)}(m;c)=\frac1N\sum_{s=0}^{p-1}m_sH_{s+j}(c),
 \qquad
 \Psi_*^{(N)}(m;c)=\min_j\Psi_j^{(N)}(m;c).
 \label{eq:Psi-definition}
\end{equation}
For $N=d$, this is exactly $D_1(m;c)$.

\begin{theorem}[Exact constant-convolution criterion]
\label{thm:saturation}
For every $N\geq1$ and $m\in\Comp_p(N)$,
\begin{equation}
 \Psi_*^{(N)}(m;c)\leq\bar c.
 \label{eq:spectral-mean-bound}
\end{equation}
The following are equivalent:
\begin{enumerate}[label=(\alph*)]
  \item $\Psi_*^{(N)}(m;c)=\bar c$;
  \item $\sum_s m_sc_{s+k}$ is independent of $k\in\Z/p\Z$;
  \item $\prod_s a_{s+k}^{m_s}$ is independent of
  $k\in\Z/p\Z$.
\end{enumerate}
\end{theorem}

\begin{proof}
Define $b_k=N^{-1}\sum_s m_sc_{s+k}$.  By linearity and covariance,
$H_j(b)=\Psi_j^{(N)}(m;c)$.  Moreover,
\[
 \frac1p\sum_kb_k
 =\frac{1}{Np}\sum_{s,k}m_sc_{s+k}
 =\bar c,
\]
and \cref{lem:H-invertible} shows that the mean of the $p$ values $H_j(b)$
is also $\bar c$.  Their minimum is at most their mean, proving
\eqref{eq:spectral-mean-bound}.  Equality holds exactly when all $H_j(b)$
equal the mean.  By \eqref{eq:H-constant-iff}, this is equivalent to $b$
being constant, which is (b).  Finally,
\[
 \sum_sm_sc_{s+k}=\log\Bigl(\prod_sa_{s+k}^{m_s}\Bigr),
\]
and injectivity of the logarithm gives the equivalence with (c).
\end{proof}

The integer form (c) is useful both conceptually and computationally: exact
saturation can be certified by shifted products, with no tolerance applied to
linear combinations of logarithms.

\subsection{What divisibility does and does not say}

For $q\in\Z/p\Z$, use the discrete Fourier convention
\begin{equation}
 \widehat c(q)=\sum_{j=0}^{p-1}c_j
 \exp\!\left(-\frac{2\pi\mathrm i qj}{p}\right).
 \label{eq:DFT}
\end{equation}

\begin{theorem}[Divisibility under full Fourier support]
\label{thm:fourier-divisibility}
Let $N\geq1$.
\begin{enumerate}[label=(\roman*)]
  \item If $p\mid N$, the uniform composition $m_s=N/p$ saturates
  \eqref{eq:spectral-mean-bound}.
  \item Suppose $\widehat c(q)\neq0$ for every $q=1,\ldots,p-1$.  Then a
  saturating composition in $\Comp_p(N)$ exists if and only if $p\mid N$,
  and every saturating composition is uniform.
\end{enumerate}
In particular, the one-level sufficient condition is $p\mid d$.  Its
necessity requires the hypothesis in (ii).
\end{theorem}

\begin{proof}
If $m_s=N/p$, then
$\sum_s m_sc_{s+k}=(N/p)\sum_sc_s$ for every $k$, so (i) follows from
\cref{thm:saturation}.

For (ii), take the Fourier transform of
$b_k=N^{-1}\sum_sm_sc_{s+k}$.  With the convention
\eqref{eq:DFT}, direct reindexing gives
\begin{equation}
 \widehat b(q)=N^{-1}\widehat m(-q)\widehat c(q).
 \label{eq:fourier-product}
\end{equation}
If $m$ saturates, $b$ is constant by \cref{thm:saturation}; hence
$\widehat b(q)=0$ for every nonzero mode.  Full nonzero Fourier support of
$c$ then forces all nonzero modes of $m$ to vanish.  Fourier inversion makes
$m$ constant.  Its integer entries sum to $N$, so this is possible exactly
when $p\mid N$, and then $m_s=N/p$.  The converse is (i).
\end{proof}

\begin{example}[A mandatory nondivisible witness]
\label{ex:nondivisible}
Take
\[
 p=4,\qquad d=N=2,\qquad a=(2,3,2,3),\qquad m=(1,1,0,0).
\]
For all four shifts $k$,
\[
 \prod_{s=0}^3a_{s+k}^{m_s}=6.
\]
Thus $m$ saturates by \cref{thm:saturation}, even though $4\nmid2$.  The
phase vector has period two, so its length-four Fourier transform has a
missing nonzero mode.  This example rules out any unconditional
``saturation if and only if $p\mid d$'' statement.
\end{example}

The distinction is between arithmetic availability and algebraic degeneracy.
Divisibility makes the uniform integer allocation available for every phase
profile.  When the profile loses Fourier modes, a nonuniform allocation can
be invisible to the circular convolution and can saturate for a different
reason.
